\subsection{网络最大流}

\subsubsection{Edmonds-Karp 算法}

每次用 BFS 找一条增广路并增广，复杂度 $\mathcal{O}(n m^2)$。

\begin{lstlisting}
using i64 = long long;
constexpr i64 INF = 1e18;
struct EdmondsKarp {
    int n;
    struct Edge {
        int to = 0,nxt = 0;
        i64 f = 0,c = 0;
    };
    std::vector<int> head,fr;
    std::vector<Edge> edge;
    EdmondsKarp(int _n):n(_n),head(n+1),edge(2),fr(n+1){}
    void addEdge(int u,int v,i64 f,i64 c){
        edge.push_back({v,head[u],f,c});
        head[u] = edge.size() - 1;
    }
    void connect(int u,int v,i64 w){
        addEdge(u,v,w,w);
        addEdge(v,u,0,w);
    }
    i64 work(int s,int t){
        i64 ans = 0;
        while (1) {
            std::queue<int> que;
            std::fill(fr.begin(),fr.end(),-1);
            que.push(s);
            fr[s] = 0;
            while (!que.empty()){
                int u = que.front();
                que.pop();
                for (int e = head[u];e;e = edge[e].nxt){
                    int v = edge[e].to;
                    i64 w = edge[e].f;
                    if (fr[v] != -1 || w == 0) continue;
                    fr[v] = e^1;
                    que.push(v);
                }
            }
            if (fr[t] == -1) break;
            i64 mn = INF;
            for (int u = t;u != s;u = edge[fr[u]].to){
                mn = std::min(mn,edge[fr[u] ^ 1].f);
            }
            for (int u = t;u != s;u = edge[fr[u]].to){
                edge[fr[u]].f += mn;
                edge[fr[u] ^ 1].f -= mn;
            }
            ans += mn;
        }
        return ans;
    }
};
\end{lstlisting}

\subsubsection{Dinic 算法}

分层图 + 多路增广，复杂度 $\mathcal{O}(n^2 m)$。对于容量均为 $1$ 的图总复杂度为 $O(m\cdot \min(m^{\frac{1}{2}},n^{\frac{2}{3}}))$；二分图最大匹配（无权）为 $O(n^{\frac{1}{2}}\cdot m)$。

\begin{lstlisting}
using i64 = long long;
constexpr i64 INF = 1e18;
struct Dinic {
    int n;
    std::vector<int> dis,now,head;
    struct Edge {
        int to = 0,nxt = 0;
        i64 f = 0,c = 0;
    };
    std::vector<Edge> edge;
    Dinic(int _n):n(_n),dis(_n+1),now(n+1),head(n+1),edge(2) {}
    void addEdge(int u,int v,i64 f,i64 c){
        edge.push_back({v,head[u],f,c});
        head[u] = edge.size() - 1;
    }
    void connect(int u,int v,i64 w){
        addEdge(u,v,w,w);
        addEdge(v,u,0,w);
    }
    bool level_graph(int s,int t){
        std::queue<int> que;
        std::fill(dis.begin(),dis.end(),-1);
        dis[s] = 0,que.push(s);
        now[s] = head[s];
        while (!que.empty()) {
            auto u = que.front();
            que.pop();
            for (int e = head[u];e;e = edge[e].nxt){
                int v = edge[e].to;
                i64 w = edge[e].f;
                if (dis[v] != -1 || w == 0) continue;
                dis[v] = dis[u] + 1;
                now[v] = head[v];
                que.push(v);
                if (v == t) return 1;
            }
        }
        return 0;
    }
    i64 blocking_flows(int u,i64 cap,const int& t){
        if (u == t) return cap;
        i64 sum = 0;
        for (int e = now[u];e;e = edge[e].nxt){
            int v = edge[e].to;
            i64 w = edge[e].f;
            if (!cap) break;
            now[u] = e;
            if (dis[v] != dis[u] + 1 || w <= 0) continue;
            i64 k = blocking_flows(v,std::min(cap,w),t);
            if (!k) dis[v] = -1;
            edge[e].f -= k,edge[e^1].f += k;
            sum += k,cap -= k;
        }
        return sum;
    }
    i64 work(int s,int t){
        i64 sum = 0;
        while (level_graph(s,t)){
            sum += blocking_flows(s,INF,t);
        }
        return sum;
    }
};
\end{lstlisting}

\subsubsection{预留推进网络最大流(HLPP)}

在稠密图或 Dinic 容易退化的网络中计算最大流；使用最高标号选择、Gap 优化与全局重标号。

常用上界记作 $O(V^2\sqrt E)$，空间 $O(V+E)$。

\begin{lstlisting}
struct HLPP {
    using i64 = long long;
    struct Edge {
        int to;
        i64 c;
        int rev;
        Edge(int to, i64 c, int rev) : to(to), c(c), rev(rev) {}
    };
    int n, s, t;
    int maxh, maxgaph, workcnt;
    std::vector<std::vector<Edge>> vec;
    std::vector<i64> ov;
    std::vector<int> h, cur;
    std::vector<int> ovList, ovNxt;
    std::vector<int> gap, gapPrv, gapNxt;

    HLPP(int n)
        : n(n), maxh(0), maxgaph(0), workcnt(0), vec(n + 1),
            ov(n + 1), h(n + 1), cur(n + 1), ovList(n + 1, -1),
            ovNxt(n + 1, -1), gap(n + 1, -1), gapPrv(n + 1, -1),
            gapNxt(n + 1, -1) {}

    void connect(int u, int v, i64 c) {
        vec[u].push_back(Edge(v, c, vec[v].size()));
        vec[v].push_back(Edge(u, 0, vec[u].size() - 1));
    }

    i64 work(int s_, int t_) {
        s = s_, t = t_;
        globalRelabel();
        for (auto &e : vec[s]) {
            if (e.c) pushFlow(s, e, e.c);
        }
        for (; maxh >= 0; --maxh) {
            while (~ovList[maxh]) {
                int x = ovList[maxh];
                ovList[maxh] = ovNxt[x];
                discharge(x);
                if (workcnt > (n << 2))
                    globalRelabel();
            }
        }
        return ov[t];
    }

private:
    void discharge(int x) {
        int nh = n, sz = vec[x].size();
        for (int i = cur[x]; i < sz; ++i) {
            auto &e = vec[x][i];
            if (e.c > 0) {
                if (h[x] == h[e.to] + 1) {
                    pushFlow(x, e, std::min(ov[x], e.c));
                    if (ov[x] == 0) {
                        cur[x] = i;
                        return;
                    }
                } else {
                    nh = std::min(nh, h[e.to] + 1);
                }
            }
        }
        for (int i = 0; i < cur[x]; ++i) {
            auto &e = vec[x][i];
            if (e.c > 0)
                nh = std::min(nh, h[e.to] + 1);
        }
        cur[x] = 0;
        ++workcnt;
        int oldh = h[x], head = gap[oldh];
        if (~gapNxt[head]) {
            setHeight(x, nh);
        } else {
            for (int i = oldh; i <= maxgaph; ++i) {
                for (int j = gap[i]; ~j; ) {
                    int nxt = gapNxt[j];
                    h[j] = n;
                    gapPrv[j] = gapNxt[j] = -1;
                    j = nxt;
                }
                gap[i] = -1;
            }
            maxgaph = oldh - 1;
        }
    }

    void globalRelabel() {
        workcnt = maxh = maxgaph = 0;
        std::fill(h.begin(), h.end(), n);
        h[t] = 0;
        std::fill(gapPrv.begin(), gapPrv.end(), -1);
        std::fill(gapNxt.begin(), gapNxt.end(), -1);
        std::fill(gap.begin(), gap.end(), -1);
        std::fill(ovList.begin(), ovList.end(), -1);
        std::fill(cur.begin(), cur.end(), 0);
        std::queue<int> q;
        q.push(t);
        while (!q.empty()) {
            int x = q.front();
            q.pop();
            for (auto &e : vec[x]) {
                if (h[e.to] == n &&
                    e.to != s &&
                    vec[e.to][e.rev].c > 0) {
                    setHeight(e.to, h[x] + 1);
                    q.push(e.to);
                }
            }
        }
    }

    void setHeight(int x, int newh) {
        if (~gapPrv[x]) {
            if (gapPrv[x] == x) {
                gap[h[x]] = gapNxt[x];
                if (~gapNxt[x])
                    gapPrv[gapNxt[x]] = gapNxt[x];
            } else {
                gapNxt[gapPrv[x]] = gapNxt[x];
                if (~gapNxt[x])
                    gapPrv[gapNxt[x]] = gapPrv[x];
            }
            gapPrv[x] = gapNxt[x] = -1;
        }
        h[x] = newh;
        if (h[x] >= n)
            return;
        maxgaph = std::max(maxgaph, h[x]);
        if (ov[x] > 0) {
            maxh = std::max(maxh, h[x]);
            ovNxt[x] = ovList[h[x]];
            ovList[h[x]] = x;
        }
        gapNxt[x] = gap[h[x]];
        if (~gapNxt[x])
            gapPrv[gapNxt[x]] = x;
        gap[h[x]] = gapPrv[x] = x;
    }

    void pushFlow(int from, Edge &e, i64 flow) {
        if (!ov[e.to] && e.to != t && h[e.to] < n) {
            maxh = std::max(maxh, h[e.to]);
            ovNxt[e.to] = ovList[h[e.to]];
            ovList[h[e.to]] = e.to;
        }
        e.c -= flow;
        vec[e.to][e.rev].c += flow;
        ov[from] -= flow;
        ov[e.to] += flow;
    }
};
\end{lstlisting}